\documentclass[11pt]{article}

\usepackage[utf8]{inputenc}
\usepackage[T1]{fontenc}
\usepackage{lmodern}
\usepackage{geometry}
\usepackage{amsmath,amssymb,amsfonts,amsthm,mathtools}
\usepackage{bm}
\usepackage{enumitem}
\usepackage{microtype}
\usepackage{hyperref}
\usepackage{titlesec}
\usepackage{booktabs}
\usepackage{array}
\usepackage{setspace}
\usepackage{mathrsfs}
\usepackage{physics}

\geometry{margin=1in}
\hypersetup{colorlinks=true, linkcolor=black, urlcolor=blue, citecolor=black}
\setlist[enumerate]{topsep=0.4em,itemsep=0.35em}
\setlist[itemize]{topsep=0.3em,itemsep=0.25em}
\titleformat{\section}{\large\bfseries}{\thesection.}{0.5em}{}
\titleformat{\subsection}{\normalsize\bfseries}{\thesubsection.}{0.5em}{}
\setstretch{1.06}

\newcommand{\Aether}{\AE{}ther}
\newcommand{\Aflow}{\AE{}ther-flow}
\newcommand{\TheoryName}{\Aflow{} Interpretation of Relativity}
\newcommand{\TheoryProgram}{\Aether{} / \Aflow{} framework}
\newcommand{\Sub}{\mathcal{A}}
\newcommand{\BridgeMetric}{\mathfrak{g}}

\newtheorem{definition}{Definition}
\newtheorem{proposition}{Proposition}
\newtheorem{remark}{Remark}
\newtheorem{corollary}{Corollary}
\newtheorem{theorem}{Theorem}

\title{\textbf{The \TheoryName{}}\\[0.4em]
\large Higher-Derivative Quartic Branch Broader Admissible-Background Theorem on Uniformly Controlled Regions}
\author{}
\date{}

\begin{document}

\maketitle

\begin{abstract}
This manuscript begins the chosen theorem-route follow-on to the quartic branch admissible-background requirements note. The scope remains explicit. The note does not claim a theorem on arbitrary admissible backgrounds, a completed nonlinear covariant derivation of Einsteinian gravity from substrate variables, or nonlinear stability of the quartic branch. It asks a narrower next question: can the already-recorded region-level admissible-background package be sharpened into a genuine broader theorem on a clearly delimited class of regions before the program turns to nonlinear stability?

The answer is yes for uniformly controlled bounded-geometry regions. Let \(\mathcal U\) be an admissible exact-closure background region equipped with a quartic branch observer structure, the previously stated admissible-background package, and a finite-overlap atlas of slowly varying admissible patches with a smooth partition of unity whose derivatives are controlled by the same background scale. Then the local curved-background continuation patches together to a region-level observer-accessible reduction on \(\mathcal U\). The reduced branch still propagates matter and light through the same single metric \(g_{\mu\nu}\), produces no independent observer-accessible tensor or transverse-vector branch correction at the present scope, and leaves only a metric-mediated scalar self-energy bounded by
\[
\Pi_{\phi,\mathcal U}^{(4)}
=
\mathcal O\!\bigl(m_{S,\min}^2\sigma_0^2\,\varepsilon_{4,\mathcal U}\bigr),
\]
where \(\varepsilon_{4,\mathcal U}\ll 1\) is the corresponding uniform recovery-compatible control parameter on the claimed mode family. In that precise sense, the quartic branch now admits a broader admissible-background theorem on uniformly controlled regions. Extending that result beyond the present bounded-geometry and finite-overlap hypotheses, and confronting nonlinear stability, remain later burdens.
\end{abstract}

\section{Introduction}

The quartic branch has now passed a full gate sequence at the disciplined scope used by the active derivational program: the flat-background exact-closure screen, the bridge-metric matter-coupling gate, the exact-relativistic recovery-compatibility gate, the first local curved-background continuation, and the region-level admissible-background requirements note. That sequence leaves a real choice for the next step. One could turn immediately to nonlinear stability, or one could first sharpen the admissible-background package into the narrowest honest broader theorem supported by the current mathematics.

The present note makes that choice explicit and records its first result. The theorem route is taken first because the current quartic branch has already isolated its background-side control data, while nonlinear stability would still depend on a more fully fixed nonlinear covariant completion than the active sequence has yet justified. The next honest step is therefore to upgrade the package note into a true theorem on a clearly defined class of regions, rather than to blur a still-local continuation into an unproved global claim.

\paragraph{Framework and claim boundary.}
\TheoryProgram{} denotes the broader \Aether{} / \Aflow{} framework built on the claims that the \Aether{} is the underlying four-dimensional substrate of reality, the \Aflow{} is its intrinsic ordered motion, observed three-dimensional space is the local experiential slice of that deeper substrate, S-time is the experienced order of change arising from matter, light, and the \Aflow{}, and observed expansion is the three-dimensional appearance of deeper four-dimensional ordered motion. Within that framework, \TheoryName{} denotes the exact-closure theory in which the effective gravitational dynamics are adopted to be exactly Einsteinian with universal matter coupling. In the active sequence, \emph{adoption} means use of that established relativistic dynamics without claiming substrate derivation, while \emph{derivation} is reserved for a first-principles recovery from explicit substrate variables.

The present note stays within that boundary. It does not claim an arbitrary-background theorem. It proves only the broader theorem now honestly supported by the local continuation and admissibility package already on record.

\section{Uniformly Controlled Regions}

\begin{definition}[Uniformly controlled quartic admissible region]
Let \(\mathcal U\) be an admissible exact-closure background region equipped with a quartic branch observer structure and satisfying the quartic admissible-background package of the preceding note. A \emph{uniformly controlled quartic admissible region} is such a region together with:
\begin{enumerate}
    \item a finite-overlap open cover \(\{\mathcal U_\alpha\}_{\alpha\in A}\) of \(\mathcal U\) by slowly varying admissible patches
    \item a smooth partition of unity \(\{\zeta_\alpha\}_{\alpha\in A}\) subordinate to that cover
    \item a background length scale \(L_{\mathcal R}\) on \(\mathcal U\) satisfying
    \begin{equation}
    \sup_{\mathcal U}\mathcal R \lesssim L_{\mathcal R}^{-2},
    \qquad
    \sup_{\mathcal U}|\nabla \mathcal R| \lesssim L_{\mathcal R}^{-3},
    \label{eq:LR}
    \end{equation}
    and partition bounds
    \begin{equation}
    |D\zeta_\alpha|\lesssim L_{\mathcal R}^{-1},
    \qquad
    |D^2\zeta_\alpha|\lesssim L_{\mathcal R}^{-2}
    \label{eq:zetabounds}
    \end{equation}
    uniformly in \(\alpha\)
    \item a finite overlap number \(N_{\mathrm{ov}}\) such that every point of \(\mathcal U\) lies in at most \(N_{\mathrm{ov}}\) patches.
\end{enumerate}
\end{definition}

\begin{definition}[Uniform region control parameter]
For a uniformly controlled quartic admissible region \(\mathcal U\), define the region-wise recovery-compatible control parameter on the claimed observer-accessible mode family by
\begin{equation}
\varepsilon_{4,\mathcal U}
\equiv
\sup_{\mathcal U}
\frac{1}{m_S^2\Omega^2}
\left[
\frac{\lambda_4}{M_*^2}k_\perp^4
+ \mathcal O\!\left(\frac{\mathcal R\,k_\perp^2}{M_*^2}\right)
+ \mathcal O\!\left(\frac{\mathcal R^2}{M_*^2}\right)
\right].
\label{eq:epsU}
\end{equation}
The region is said to lie in the uniform recovery-compatible regime when \(\varepsilon_{4,\mathcal U}\ll 1\).
\end{definition}

\begin{remark}
Definition 2 strengthens the earlier admissible-background package only by adding the finite-overlap patching data needed to upgrade patchwise control into a theorem on the whole region. No arbitrary-background claim is hidden in this step.
\end{remark}

\section{Uniform Patching of the Local Continuation}

\begin{proposition}[Uniform local patch estimates]
Let \(\mathcal U\) be a uniformly controlled quartic admissible region in the regime \(\varepsilon_{4,\mathcal U}\ll 1\). Then each patch \(\mathcal U_\alpha\) obeys the previously established local continuation theorem with constants that may be chosen uniformly in \(\alpha\). In particular, on each patch one has
\begin{equation}
\mathcal B_{4,\alpha}^{\mathrm{cov}}
=
\frac{\lambda_4}{M_*^2}(-D^2)^2
+ \mathcal O\!\left(\frac{\mathcal R\,D^2}{M_*^2}\right)
+ \mathcal O\!\left(\frac{\nabla\mathcal R\,D}{M_*^3}\right)
+ \mathcal O\!\left(\frac{\mathcal R^2}{M_*^2}\right),
\label{eq:Balpha}
\end{equation}
and
\begin{equation}
\Pi_{\phi,\alpha}^{(4)}
=
\mathcal O\!\bigl(m_{S,\min}^2\sigma_0^2\,\varepsilon_{4,\mathcal U}\bigr).
\label{eq:Pialpha}
\end{equation}
\end{proposition}

\begin{proof}
The preceding admissible-background requirements note already showed that the admissible-background package supplies the exact hypotheses needed to reduce every sufficiently small neighborhood of \(\mathcal U\) to the local curved-background continuation theorem. The additional data of Definition 1 ensure that the same curvature scale, partition regularity scale, and overlap number may be chosen uniformly across the cover. Therefore the local operator form and the local self-energy estimate inherit constants independent of \(\alpha\), yielding \eqref{eq:Balpha} and \eqref{eq:Pialpha}.
\end{proof}

\begin{proposition}[Glued quartic operator on the region]
Under the hypotheses of Proposition 1, the local quartic operators patch to a region-level reduced scalar operator
\begin{equation}
\mathcal B_{4,\mathcal U}
=
\sum_{\alpha\in A}\zeta_\alpha\,\mathcal B_{4,\alpha}^{\mathrm{cov}}\,\zeta_\alpha
+ \delta\mathcal B_{4,\mathrm{ov}},
\label{eq:BU}
\end{equation}
where the overlap remainder satisfies the same schematic bound
\begin{equation}
\delta\mathcal B_{4,\mathrm{ov}}
=
\mathcal O\!\left(\frac{\mathcal R\,D^2}{M_*^2}\right)
+ \mathcal O\!\left(\frac{\nabla\mathcal R\,D}{M_*^3}\right)
+ \mathcal O\!\left(\frac{\mathcal R^2}{M_*^2}\right)
\label{eq:Boverlap}
\end{equation}
on \(\mathcal U\), with constants depending only on the uniform control data and \(N_{\mathrm{ov}}\).
\end{proposition}

\begin{proof}
Insert the partition of unity \(1=\sum_\alpha \zeta_\alpha^2\) into the local reduced scalar sector and commute the derivatives in \(\mathcal B_{4,\alpha}^{\mathrm{cov}}\) past the cutoff functions. The resulting commutators produce derivatives of \(\zeta_\alpha\), which are bounded by \eqref{eq:zetabounds}, and at most the same differential order already present in \(\mathcal B_{4,\alpha}^{\mathrm{cov}}\). Using \eqref{eq:LR}, the partition bounds, and finite overlap, those commutator terms are of the same schematic sizes as the curvature-suppressed terms already kept in \eqref{eq:Balpha}. Summing over the finite-overlap atlas therefore yields \eqref{eq:BU} and \eqref{eq:Boverlap}.
\end{proof}

\begin{remark}
Proposition 2 is the actual theorem-route upgrade over the earlier package note. The earlier note identified what a broader theorem would need to control. Proposition 2 shows that, on a defined class of uniformly controlled regions, the patchwise local theorem really does glue into a region-level reduced operator.
\end{remark}

\section{Region-Level Observer-Accessible Reduction}

\begin{proposition}[Uniform region-level self-energy bound]
On a uniformly controlled quartic admissible region in the regime \(\varepsilon_{4,\mathcal U}\ll 1\), the reduced observer-accessible branch correction remains a metric-mediated scalar self-energy satisfying
\begin{equation}
\Pi_{\phi,\mathcal U}^{(4)}
=
\mathcal O\!\bigl(m_{S,\min}^2\sigma_0^2\,\varepsilon_{4,\mathcal U}\bigr).
\label{eq:PiU}
\end{equation}
No independent observer-accessible tensor or transverse-vector branch correction is generated at the same scope.
\end{proposition}

\begin{proof}
The local continuation theorem and Proposition 1 show that each patch supports only the already-known metric-mediated scalar self-energy, with no tensor or transverse-vector branch correction at the present scope. Proposition 2 patches the scalar operator into the region-level operator \(\mathcal B_{4,\mathcal U}\) while preserving the same derivative and curvature hierarchy. The corresponding Gaussian elimination therefore yields the same reduced scalar structure region-wise, and the bound \eqref{eq:PiU} follows from the uniform control parameter \(\varepsilon_{4,\mathcal U}\), the uniform lower bound \(m_S^2\ge m_{S,\min}^2\), and finite overlap.
\end{proof}

\begin{corollary}[Operational relativistic sector on uniformly controlled regions]
On a uniformly controlled quartic admissible region in the regime \(\varepsilon_{4,\mathcal U}\ll 1\), the observer-accessible operational sector remains that of the same single effective metric \(g_{\mu\nu}\): one null cone, one proper-time rule, and the ordinary local SR limit on sufficiently small neighborhoods of every point of \(\mathcal U\).
\end{corollary}

\begin{proof}
The universal matter route is unchanged, and Proposition 3 shows that the only surviving branch correction remains a uniformly subleading metric-mediated scalar self-energy. Hence the operational geometry seen by matter and light is still governed by the same effective metric throughout the claimed region.
\end{proof}

\section{Broader Admissible-Background Theorem}

\begin{theorem}[Broader admissible-background theorem on uniformly controlled regions]
Let \(\mathcal U\) be a uniformly controlled quartic admissible region for the higher-derivative quartic branch, and assume the uniform recovery-compatible regime \(\varepsilon_{4,\mathcal U}\ll 1\). Then, at the present observer-accessible quadratic and WKB-style scope, the quartic branch admits a broader admissible-background continuation on \(\mathcal U\) with the following properties.
\begin{enumerate}
    \item The observer-accessible matter route remains the single metric route \(S_{\mathrm{matter}}[g,\psi]\).
    \item The reduced scalar auxiliary operator is globally represented on \(\mathcal U\) by \eqref{eq:BU}, with principal quartic part \((\lambda_4/M_*^2)(-D^2)^2\) and only curvature- and overlap-suppressed remainders of the schematic form \eqref{eq:Boverlap}.
    \item The only surviving observer-accessible branch correction remains the metric-mediated scalar self-energy \(\Pi_{\phi,\mathcal U}^{(4)}\) satisfying the bound \eqref{eq:PiU}.
    \item No independent observer-accessible tensor channel, transverse-vector channel, second null cone, or second proper-time rule is generated at the same scope.
\end{enumerate}
Accordingly, the quartic branch preserves the active exact-relativistic benchmark on uniformly controlled admissible regions of this type.
\end{theorem}

\begin{proof}
Item 1 is inherited from the universal bridge-metric matter route already fixed by the matter-coupling gate and retained in the admissible-background package. Item 2 is Proposition 2. Item 3 is Proposition 3. Item 4 follows from Proposition 3 and Corollary 1, together with the fact that the quartic branch does not modify the observer map or introduce a second matter metric. These four items are exactly the region-level conclusions needed for a broader admissible-background theorem at the present disciplined scope.
\end{proof}

\begin{remark}
Theorem 1 is broader than the earlier one-patch continuation and stronger than the earlier package note. It is still not a theorem on arbitrary admissible backgrounds, and it is still not a nonlinear stability result.
\end{remark}

\section{What This Note Establishes and What It Does Not}

The present note establishes the following.
\begin{enumerate}
    \item It records the explicit decision to push the admissible-background package toward a theorem before turning to nonlinear stability.
    \item It identifies a concrete class of uniformly controlled admissible regions on which the local continuation theorem can be patched into a genuine region-level result.
    \item It proves that, on those regions, the quartic branch remains single-metric and leaves only a uniformly subleading metric-mediated scalar self-energy at the present scope.
    \item It upgrades the earlier region-level package note into a broader admissible-background theorem on that delimited class of regions.
\end{enumerate}

It does not establish the following.
\begin{enumerate}
    \item It does not prove a theorem on arbitrary admissible backgrounds.
    \item It does not control the ultrasoft static corner at the same scope as the non-ultrasoft regime.
    \item It does not fix a unique full nonlinear covariant completion of the quartic branch.
    \item It does not solve nonlinear stability.
    \item It does not complete a first-principles substrate derivation of Einsteinian gravity.
\end{enumerate}

\section{Conclusion}

The repository has now moved one disciplined step beyond the admissible-background package note. The quartic branch no longer has only a local continuation and a list of region-level requirements. On uniformly controlled bounded-geometry regions, those ingredients now patch into an honest broader admissible-background theorem.

This remains a bounded theorem, not a declaration of final success. The next open derivational burden is narrower again: determine whether the present theorem can be extended beyond its finite-overlap bounded-geometry hypotheses or whether a disciplined obstruction appears there. Nonlinear stability remains the later burden, to be confronted only after the admissible-background reach of the quartic branch has been fixed more sharply.

\end{document}
